\section{Fields}

\begin{exercise}[Hertz's waves]
\textbf{Experiment.}
Hertz produced and detected electromagnetic waves whose reflection,
interference, and propagation speed agreed with Maxwell's theory
\cite{hertz}.  On oriented spacetime, represent the electromagnetic field by
a two-form \(F\) and current by a three-form \(J\).

Write Maxwell's equations using the exterior derivative and the metric's
Hodge operator, without splitting space from time.  Show that charge
conservation follows from applying the exterior derivative twice.  In a
source-free contractible region, derive a wave equation for \(F\).  Explain
which part of the construction needs a metric, which needs only orientation,
and why the predicted propagation speed is independent of the observer's
coordinates.
\end{exercise}

\begin{exercise}[A charge bends but does not speed up]
\textbf{Experiment.}
In a uniform magnetic field, a nonrelativistic charged particle follows a
circle at constant speed.  Let \(Q\) be physical space, let \(F\) be the
magnetic two-form on \(Q\), and let the kinetic energy supply an identification
between velocity and momentum.

Construct the twisted symplectic form
\[
        \Omega_F=\Omega_{\mathrm{can}}+q\,\pi^*F
\]
on \(T^*Q\).  Define the motion by
\(\iota_{X_H}\Omega_F=dH\) and recover the Lorentz force after choosing
Euclidean coordinates only for the final calculation.  Show directly from
antisymmetry that \(dH(X_H)=0\), explaining the constant speed.  Determine
what condition on \(F\) is forced by \(d\Omega_F=0\), and relate its failure to
the hypothetical observation of magnetic charge.
\end{exercise}

\begin{exercise}[Phase volume does not compress]
\textbf{Experiment.}
An initially small ensemble of nearly isolated classical systems stretches
and folds in phase space but does not change its fine-grained phase volume.
Let \((M,\Omega)\) be a finite-dimensional symplectic phase space and let
\(\varphi_t\) be the flow generated by \(H:M\to\R\).

Use Cartan's formula to prove
\(\mathcal L_{X_H}\Omega=0\), then prove invariance of
\(\Omega^n/n!\).  Explain how this accounts for the observed incompressibility
without choosing canonical coordinates.  Contrast the conclusion with a
damped oscillator, and identify precisely which structural assumption the
damped dynamics violates.
\end{exercise}

\begin{exercise}[LIGO's correlated motion]
\textbf{Experiment.}
LIGO observed a time-dependent differential change in the lengths of two
perpendicular arms consistent with a passing gravitational wave
\cite{ligo}.  Model freely falling mirrors by nearby geodesics with separation
field \(J\) and tangent \(u\).

Derive the geodesic-deviation equation
\[
        \nabla_u\nabla_uJ+R(J,u)u=0.
\]
Show that the relative acceleration is curvature and hence cannot be removed
throughout the apparatus by a change of coordinates.  In the weak-field
limit, explain why a transverse traceless disturbance lengthens one arm while
shortening the other.  Identify the measured quantity as a comparison of
proper times or optical path lengths, not as the value of a coordinate
component.
\end{exercise}

\begin{exercise}[The charged weak bosons]
\textbf{Experiment.}
Collider events revealed massive charged \(W^\pm\) bosons and a neutral
\(Z^0\) boson with the weak-interaction decay patterns predicted by the
electroweak model \cite{ua1-w,ua1-z}.  Let a gauge group \(G\) act on a vector
bundle \(E\to M\), let \(A\) be a connection, and let
\(F_A\) be its curvature.

Show how a local change of frame transforms \(A\) inhomogeneously but
transforms \(F_A\) by conjugation.  Deduce that an invariant pairing on the
Lie algebra gives a frame-independent Yang--Mills action.  Let a nonzero
scalar field select a stabilizer subgroup \(H\subset G\).  Explain why
directions tangent to \(G/H\) acquire mass while the unbroken electromagnetic
direction remains massless.  Relate the resulting charged and neutral
directions to the observed particle pattern without treating a gauge choice
as an observable.
\end{exercise}

\begin{exercise}[Faraday induction]
\textbf{Experiment.}
A changing magnetic flux through a conducting loop produces an electromotive
force around the loop.  Let \(F\) be the spacetime electromagnetic two-form
and let the loop sweep out a cylinder between two times.

Apply Stokes' theorem to the cylinder and derive Faraday's law as a comparison
of fluxes and line integrals.  Explain why the result remains meaningful for a
moving loop and why no coordinate-dependent electric field components are
fundamental to the statement.  Determine the sign from the orientation and
explain Lenz's law as the experimentally observed energetic response of the
circuit.
\end{exercise}

\begin{exercise}[Foucault's pendulum]
\textbf{Experiment.}
The oscillation plane of a long pendulum rotates over a day by an amount which
depends on latitude.  Treat the slowly transported oscillation direction as a
tangent vector on the spherical Earth.

Parallel transport the direction around the latitude circle and relate its
net rotation to the enclosed Gaussian curvature.  Recover the observed
latitude dependence.  Explain why the angle is holonomy rather than a failure
of the vector to remain fixed, and separate this geometric effect from local
friction and from the choice of longitude coordinates.
\end{exercise}

\begin{exercise}[Gravitational redshift]
\textbf{Experiment.}
Identical clocks placed at different heights in a stationary gravitational
field accumulate different proper times, as confirmed by clock-comparison
experiments.  Let the stationary spacetime possess a timelike symmetry field
\(T\).

Show that the conserved photon energy along a null geodesic is obtained by
pairing its tangent with \(T\).  Relate the frequencies measured by stationary
observers to the norms of their four-velocities and derive the redshift ratio.
Explain why the prediction compares proper times and survives all coordinate
changes.  Recover the leading weak-field height dependence.
\end{exercise}

\begin{exercise}[Light is gravitationally lensed]
\textbf{Experiment.}
Light from a distant source can reach an observer along multiple paths around
a massive object, producing several images and measurable time delays.
Treat light rays as null geodesics of spacetime.

Use the variational characterization of null paths to explain the occurrence
of stationary optical paths.  Show why curvature can produce more than one
such path with the same endpoints.  Express image distortion through the
differential of the endpoint map and identify conjugate points as its
singularities.  Explain why image positions and relative delays are
observables although any drawn coordinate grid is not.
\end{exercise}

\begin{exercise}[The Sagnac time difference]
\textbf{Experiment.}
Two light beams sent in opposite directions around a rotating loop return at
different times; ring-laser gyroscopes measure the resulting phase shift.
Model the apparatus by a closed world tube and the light beams by its two null
directions.

Derive the proper-time difference measured by the emitter and show that its
leading term is proportional to the loop's oriented area and angular
velocity.  Express the result as a circulation associated with the observer
field.  Explain why the phase detects rotation of the apparatus and not an
arbitrary rotation of spatial coordinates.
\end{exercise}
